% HD-F1-T41. GPT-6 Astra Pro, /root/transport; 22 November 2026 UTC.
\section[HD-F1-T41: the remaining orthogonal rank bounds]{HD-F1-T41: an effective rank bound for the remaining orthogonal sources}
\label{sec:f1-remaining-orthogonal-rank-cutoff}

Put \(L_0=3\cdot10^{30}-3\). Every original remaining
T2-OTHER source with projective socle
\(P\Omega_{2r+1}(q)\), for odd \(q\), or
\(P\Omega^-_{2r+2}(q)\), in either characteristic, satisfies
\[
             \boxed{\qquad r\le4L_0^2+8<10^{62}.\qquad}
 \tag{F1T41.1}
\]
Here \(r\) is the Witt index; the retained field bound is
\(q\le27\). The odd-dimensional branch is explicitly
restricted to odd characteristic. The actual odd-dimensional
quadratic quotient in characteristic two is the polar-radical
isomorphism to the symplectic group proved in T36, and has
the T39 bound through that same map.

The new ingredient is the actual anisotropic coefficient
space in a partial Fourier model. Its Heisenberg operators
both force the required lower-level code conditions and
prove the full-weight stabilizer's irreducibility. No
restriction from a larger split orthogonal representation
is substituted for an original constituent.

\paragraph{1. Exact natural models and all original automorphisms.}
Retain the original \(G\subset GL_K(V)\),
\(m=\dim_KV\ge27\), scalar centre \(Z\), perfect
layer \(L\), and T35--T36 natural surjection
\[
 \begin{aligned}
 \beta:H=\Omega(A,Q)\longrightarrow L,
 &\qquad A=X\oplus A_0\oplus X^*,\quad \dim X=r,\\
 Q(x,z,y)=x^ty+Q_0(z),
 &\qquad a:=\dim A_0\in\{1,2\}.
 \end{aligned}
 \tag{F1T41.2}
\]
The polar form \(B_0\) of \(Q_0\) is nondegenerate
in exactly the characteristic scopes just stated.
Write \(N=2r+a\). Suppose \(r\ge4L_0^2+9\).
This excludes all covering exceptions and triality ranks.

The following coordinate comparisons are actual maps.
For odd-dimensional \(A\), write \(Q_0(z)=c z^2\).
The map \((x,z,y)\mapsto(x,z,c^{-1}y)\) pulls back
\(x^ty+z^2\) to \(c^{-1}Q\). It conjugates the
actual isometry and derived groups, since multiplication
of a form by a nonzero scalar leaves its isometries
unchanged. Transport \(\beta\) by its inverse conjugation.
For a nonsplit plane \(Q_0(x,y)=b x^2+dxy+e y^2\),
anisotropy makes
\(T^2-(d/b)T+e/b\) irreducible. If \(t\) is a root
in \(\mathbf F_{q^2}\) and \(z_0^{q+1}=b\), then
\[
 (x,y)\longmapsto z_0(x+t y),\qquad
 N_{q^2/q}\bigl(z_0(x+t y)\bigr)=Q_0(x,y)
 \tag{F1T41.3}
\]
is the required actual isometry to the norm plane.
Thus use \(Q_0=z^2\) in the first branch and
\(Q_0(z)=z^{q+1}\) in the second, with the displayed
transported \(\beta\).

Steinberg's finite-field Theorem 30 gives the inner,
diagonal and field actions for type \(B_r\); there is
no diagram action in this range. Given simple-root
multipliers \(c_1,\ldots,c_r\), put
\(a_r=c_r\), \(a_i=c_i a_{i+1}\).
The actual isometry
\(\operatorname{diag}(a_i,1,a_i^{-1})\)
has multipliers \(a_i/a_{i+1}\) and \(a_r\) on
the simple roots \(\varepsilon_i-\varepsilon_{i+1}\)
and \(\varepsilon_r\). For the last multiplier,
one can check the upper short-root formula obtained
from (F1T41.11) below by interchanging \(x_i,y_i\).
Field actions in these coordinates are entrywise Frobenius.

For \({}^2D_{r+1}\), Steinberg Theorem 36 gives inner,
diagonal and field actions. Its supplied proof is a
sketch, with its exact access scope recorded separately.
The diagonal representatives can again be exhibited.
For \(c_1,\ldots,c_{r-1}\in\mathbf F_q^\times\)
and \(t\in\mathbf F_{q^2}^\times\), set
\[
 a_r=1,\quad a_i=\prod_{j=i}^{r-1}c_j,
 \quad c=t^{-1},\quad\lambda=c^{q+1},
 \quad d(x,z,y)=(\operatorname{diag}(a_i)x,c z,
                         \lambda\operatorname{diag}(a_i^{-1})y).
 \tag{F1T41.4}
\]
This scales \(Q\) by \(\lambda\). Over \(\mathbf F_{q^2}\),
the two coordinates of the norm plane have multipliers
\(c,c^q\); the two terminal-root multipliers are
therefore \(t,t^q\), and the other multipliers are
the prescribed \(c_i\). These are precisely the
twist-compatible diagonal characters. The field action
is \((x,z,y)\mapsto(x^{p^e},z^{p^e},y^{p^e})\),
where \(z\) is in \(\mathbf F_{q^2}\).
Its \(q\)-power interchanges the norm-plane coordinates;
this includes the terminal-node action within the
twisted field factor, without omitting it.

Consequently every original automorphism has an actual
semilinear similitude \(T_g\) of this same natural space
with
\[
 Q(T_gv)=\lambda_g Q(v)^{p^{e_g}},\qquad
 \alpha_g(h)=T_g hT_g^{-1},\qquad
                         \beta\alpha_g=c_g\beta.
 \tag{F1T41.5}
\]
The last equality follows from the actual T36 lift:
two lifts of the same projective action differ by a
homomorphism from the perfect \(H\) to its centre,
which is trivial. This also retains the original natural
central kernel and every outer action.

Take a nonidentity difference-root element in the
hyperbolic block, as displayed in T40. It has natural
moving rank two and nonidentity image in the simple
projective socle. Its \(H\)-conjugates generate that
socle, since their generated subgroup is nontrivial and
normal. For every \(g\notin Z\), one conjugate \(u\)
therefore has \([g,\beta(u)]\) projectively nontrivial;
otherwise the original projective \(g\) centralizes
the socle in its actual almost simple group. Put
\[
 h=\alpha_g(u)u^{-1},\qquad
 \beta(h)=[g,\beta(u)],\qquad
                    2\le\operatorname{rank}_A(h-I)\le4.
 \tag{F1T41.6}
\]
The upper bound follows from (F1T41.5). The nonzero
even rank is exactly GLT Lemma 3.3's parity input,
with its GT Lemma 5.8(ii) dependency in even characteristic,
as already scoped in T40. No inertia operator is rescaled.

\paragraph{2. Genuine smaller symplectic modules, including levels one and two.}
Fix a nontrivial additive character \(\psi\) of
\(\mathbf F_q\), and put
\[
 U=\mathbf F_q^k\oplus\mathbf F_q^k,
 \quad b_U((x,z),(x',z'))=z\cdot x'-x\cdot z',
 \quad S=\operatorname{Sp}(U),\quad
 \mathcal H_0=\mathbf C[A_0^k],\quad\dim\mathcal H_0=q^{ak}.
 \tag{F1T41.7}
\]
For \(a_0\in A_0\), \(u=(x,z)\in U\), define
the actual operator
\[
 W_{a_0}(u)\delta_w=
 \psi\!\left(\sum_j z_j B_0(w_j,a_0)
                        +Q_0(a_0)x\cdot z\right)
                      \delta_{w+x a_0}.
 \tag{F1T41.8}
\]
Translations and multiplications among these operators
span the entire endomorphism algebra of \(\mathcal H_0\):
nondegeneracy of \(B_0\) gives all additive characters,
their Fourier averages give the individual diagonal
projections, and translations give every matrix unit.

There is a genuine unitary \(S\)-action \(\rho_0\) on
this actual space satisfying
\[
       \rho_0(s)W_{a_0}(u)\rho_0(s)^{-1}=W_{a_0}(su).
 \tag{F1T41.9}
\]
Here is an actual construction retaining small-rank phases.
In odd characteristic diagonalize \(Q_0\) and tensor
the genuine parameter Weil modules used in T37--T38.
In even characteristic use the genuine norm-plane
unitary-Heisenberg extension of T39. Its unitary element
\(((a_0x,a_0z),0)\) has the coordinate centre
\(Q_0(a_0)x\cdot z\) in T39's Schr\"odinger coordinates,
so its action is exactly (F1T41.8).

To avoid imposing any unproved low-rank normalization,
perform these constructions first at symplectic rank
\(k+5\). Restrict the genuine group action to its
first \(k\) coordinates. It commutes with the entire
translation/multiplication algebra of the last five
coordinates, so is an actual operator
\(\rho_0(s)\otimes I\) on
\(\mathbf C[A_0^k]\otimes\mathbf C[A_0^5]\).
The inclusion \(v\mapsto v\otimes\delta_0\) and
evaluation against \(\delta_0\) recover \(\rho_0(s)\).
They prove the group law and (F1T41.9), retaining the
actual scalar of every operator. This works for every
\(k\ge1\), including \(k=1,2\).

\paragraph{3. The actual original permutation module and its commuting action.}
GLT Lemma 3.3 and the explicit averaged-Hom intertwiner
of T40 identify the two genuine level generators on
\(H\). The T36 level bound therefore gives, for each
original constituent, an actual injection into
\(\mathcal R_k=\mathbf C[A^k]\) of minimal permutation
level \(1\le k\le L_0\). Use the partial Fourier map
in its hyperbolic \(y\)-coordinates:
\[
 e_{x,z,w}=q^{-rk/2}\sum_y\psi(\operatorname{Tr}(y^tz))
                                      \delta_{(x,w,y)}.
 \tag{F1T41.10}
\]
It is a unitary identification
\(\mathcal R_k=\mathbf C[\operatorname{Hom}(X,U)]
\otimes\mathcal H_0\), of dimension \(D=q^{Nk}\).
The rows \(F_i=(x_i,z_i)\) are the corresponding
map \(F:X\to U\). Define the genuine smaller action
by \(e_F\otimes v\mapsto e_{sF}\otimes\rho_0(s)v\).

This action commutes with the entire original \(H\),
as can be checked before the Fourier basis change.
For a smaller symplectic Levi, its operator is an actual
scalar times the simultaneous change of the \(k\)
columns in \(A^k\). For a smaller symmetric shear,
it is a scalar times multiplication by the character
of the corresponding quadratic combination of
\(Q(v_i)\) and \(B_A(v_i,v_j)\).
For a smaller one-coordinate Weyl matrix, it is a
scalar times the normalized Fourier transform with
kernel \(\psi(B_A(v_i,v_i'))\).
These assertions follow by substitution in (F1T41.10)
and (F1T41.8)--(F1T41.9); uniqueness up to the retained
scalar follows from the matrix-unit argument above.
Every one of these three operators commutes with every
isometry of \(A\), acting by its actual permutation
on \(A^k\). These generators generate \(S\), proving
the claimed commuting \(H\times S\)-action.

The large hyperbolic Levi, alternating Siegel group
and paired Fourier operators are exactly those of
T40, acting trivially on \(\mathcal H_0\). There are
also actual orthogonal unipotents, for \(a_0\in A_0\),
\[
 u_{i,a_0}:(x,z,y)\longmapsto
 \left(x,z+x_i a_0,
 y-e_i\bigl(B_0(z,a_0)+x_iQ_0(a_0)\bigr)\right).
 \tag{F1T41.11}
\]
Substitution proves that these preserve \(Q\) and
\(u_{i,a_0}u_{i,b_0}=u_{i,a_0+b_0}\).
For nonzero \(a_0\) their moving rank is two.
They belong to the actual \(H\): when \(q>2\),
conjugation by the determinant-one hyperbolic Levi
with entries \(t,t^{-1}\) at \(i\) and another
coordinate sends \(u_{i,b_0}\) to \(u_{i,t^{-1}b_0}\).
Choose \(t\ne1\) and express (F1T41.11) as this
commutator with \(b_0=(t^{-1}-1)^{-1}a_0\).
The Levi is in \(\Omega\), which is normal in the
isometry group, so the commutator is in \(\Omega\).
For \(q=2\), the actual even-rank quasideterminant
kernel from GT Lemma 5.8(ii) proves membership directly.

On (F1T41.10), direct substitution gives the exact action
\[
          u_{i,a_0}(e_F\otimes v)
                  =e_F\otimes W_{a_0}(F_i)v.
 \tag{F1T41.12}
\]
In particular these operators preserve every full
alternating Siegel weight \(F^*b_U\).

\paragraph{4. Exact lower permutation modules inside the anisotropic model.}
For an isotropic \(J\subset U\) of dimension \(j\), set
\[
 \mathcal H_J=
 \{v:W_{a_0}(u)v=v\text{ for every }a_0\in A_0, u\in J\}.
 \tag{F1T41.13}
\]
For \(J_0\) spanned by \(j\) of the \(z\)-coordinates,
the characters in (F1T41.8) show that this is exactly
the span of \(\delta_w\) with the same \(j\) entries
of \(w\) zero. A symplectic basis map \(sJ_0=J\)
and the actual \(\rho_0(s)\) in (F1T41.9) give
\[
       \mathcal H_J=\rho_0(s)\mathcal H_{J_0},
                    \qquad\dim\mathcal H_J=q^{a(k-j)}.
 \tag{F1T41.14}
\]

Define \(C_J\) as the functions supported on
\(\operatorname{Hom}(X,J^\perp)\) whose values lie in
\(\mathcal H_J\) and are constant on cosets of
\(\operatorname{Hom}(X,J)\). Its normalized
linear isometry is
\[
 i_J(e_{\bar F}\otimes v)
       =q^{-rj/2}\sum_{\Delta:X\to J}e_{F+\Delta}\otimes v,
                         \qquad v\in\mathcal H_J.
 \tag{F1T41.15}
\]
This is an actual lower-level \(H\)-module, not merely
a dimension comparison. For \(J_0\), start with the
actual inclusion \(\mathbf C[A^{k-j}]\to\mathbf C[A^k]\)
by setting its last \(j\) columns equal to zero.
Formula (F1T41.10) turns that inclusion precisely into
(F1T41.15), including its \(q^{-rj/2}\) scalar and
the zero residual coordinates in \(\mathcal H_{J_0}\).
For general \(J\), apply the actual commuting smaller
operator \(s\), and use its induced basis isometry
\(J_0^\perp/J_0\to J^\perp/J\) and \(\rho_0(s)\).
Thus
\[
       C_J\simeq\mathbf C[A^{k-j}],
                  \qquad\dim C_J=q^{N(k-j)}.
 \tag{F1T41.16}
\]

\paragraph{5. The entire lower-weight part, with the coefficient condition proved.}
For \(r\ge2k+2\) we prove
\[
 \mathcal R_k^{<2k}=\sum_{\dim J=1}C_J,
 \qquad E:=\dim\mathcal R_k^{<2k}<q^{2k+N(k-1)}.
 \tag{F1T41.17}
\]
Take an irreducible \(K\) of largest alternating
Siegel rank \(d<2k\), and a full-weight vector
\(\Phi(F)\in\mathcal H_0\). For a supported \(F\), put
\(M=\operatorname{im}F\), \(J=M\cap M^\perp\),
\(l=\dim M\), \(j=\dim J\). Then
\(d=l-j\) and \(l+j\le2k\).

The actual paired hyperbolic Fourier operators of T40
act on \(F\) and leave the coefficient space unchanged.
Their two-plane slice argument therefore applies to
each coefficient: transformed support lies in
\(\operatorname{Hom}(X_1,J^\perp)\) for
\(X_1\subset\ker F\), since otherwise the alternating
rank rises by two. Fourier orthogonality makes the
slice invariant under \(\operatorname{Hom}(X_1,J)\).
Keep a complementary subspace surjecting onto \(M\)
while composing these translations; intermediate kernel
values lie in \(J\), so the support argument persists.
It follows that
\[
 \Phi(F+\Delta)=\Phi(F)
 \quad\text{if }\Delta:X\to J
       \text{ and }F|_{\ker\Delta}\text{ surjects onto }M.
 \tag{F1T41.18}
\]

There is now an additional exact condition.
Choose an actual \(SL(X)\)-adapted basis with one
kernel coordinate zero, and hold the surjective
complement fixed. For any \(u\in J\), (F1T41.18)
equates the coefficients with that coordinate zero
and with value \(u\). Apply (F1T41.18) also to
the full-weight vector \(u_{i,a_0}\Phi\).
Equation (F1T41.12) makes its two coefficients
\(\Phi(F)\) and \(W_{a_0}(u)\Phi(F)\).
They are equal. Hence
\[
                         \Phi(F)\in\mathcal H_J.
 \tag{F1T41.19}
\]
This uses the actual large-group action; it does not
discard the anisotropic coefficient factor.

For two full-image maps in the same \(J\)-coset,
choose \(\Delta:X\to J\) uniformly. The finite
row-space bound of T40 is
\[
            2(q^l-1)(q^j-1)q^{-r}<1.
 \tag{F1T41.20}
\]
It provides a common intermediate map for both
applications of (F1T41.18). Thus their coefficient
vectors agree, also when \(q=2\).
Set \(j_0=(2k-d)/2\). On each maximal branch
\(\dim J=j_0\), one has \(M=J^\perp\).
For its surjective quotient map \(\bar F\), subtract
exactly
\[
        q^{rj_0/2}i_J(e_{\bar F}\otimes v_{J,\bar F}),
                  \qquad v_{J,\bar F}\in\mathcal H_J,
 \tag{F1T41.21}
\]
where \(v_{J,\bar F}\) is the common full-image
coefficient. Full-image lifts exist, and distinct
maximal branches do not interfere because inclusion
of their equal-dimensional orthogonal complements
forces equality of their radicals.

The remainder lies in the stable space
\(K+\sum_{\dim J=j_0}C_J\), retains its full
weight and has no maximal branch. If nonzero,
take a supported map of maximal ordinary map rank
and a two-plane in its kernel. Its slice is supported
on maps into \(J\), by maximality and equality of
the weights; (F1T41.18) makes it a nonzero constant
vector times that indicator. The actual paired
Fourier transform is nonzero at every map into
\(J^\perp\). Since
\(\dim(J^\perp/J)-d=2k-2j-d\ge2\), two added
vectors in a nondegenerate complementary plane raise
the weight rank by two, a contradiction. The
remainder is zero.

Every \(C_J\) with \(j>0\) lies in \(C_\ell\)
for \(\ell\subset J\): the support and translation
conditions imply those for \(\ell\), and
\(\mathcal H_J\subset\mathcal H_\ell\).
Conversely (F1T41.16) bounds its largest weight rank
by \(2(k-j)\). This proves the equality in
(F1T41.17), and fewer than \(q^{2k}\) lines prove
the dimension bound. As in T40, projecting from
the code sum and taking the averaged Hermitian
adjoint embeds any lower-weight constituent into
the actual previous permutation level. A minimal
level-\(k\) original constituent is therefore top rank.

\paragraph{6. The full-weight stabilizer and the exact top projector.}
Fix a surjection \(F_0:X\to U\). Its alternating
weight fibre is the actual torsor \(\{sF_0:s\in S\}\),
so its full coefficient space is
\(\mathbf C[S]\otimes\mathcal H_0\).
The smaller group acts by left translation tensor
\(\rho_0\). Use the actual unitary change of basis
\[
          e_s\otimes v\longmapsto
                            e_s\otimes\rho_0(s)^{-1}v.
 \tag{F1T41.22}
\]
It makes the smaller action exactly left regular
tensor the identity. Equation (F1T41.9) turns the
large operators (F1T41.12) into the constant
operators \(I\otimes W_{a_0}((F_0)_i)\).
Because the rows of \(F_0\) span \(U\), these
generate \(I\otimes\operatorname{End}(\mathcal H_0)\).
Indeed (F1T41.8) gives
\(W_{c a_0}(u)=W_{a_0}(cu)\), and products
add the labels up to their explicit central unit
scalars; translations and multiplications then
give all matrix units as above.

The actual \(SL(X)\) stabilizer of the weight
surjects onto \(S\) on the quotient by its radical.
Its operators become right regular translations
tensor the corresponding inverse \(\rho_0\)
operators under (F1T41.22). These are literal
operators obtained by lifting the quotient map
and acting identically on the radical; their
determinants are one. All full-rank weights
are one \(SL(X)\)-orbit, with any adapting
determinant corrected on the nonzero radical.

Write the genuine smaller isotypic decomposition as
\(\mathcal R_k=\bigoplus_\tau E_\tau\otimes D_\tau\),
with \(b=\dim E_\tau\) and \(s=|S|<q^{4k^2}\).
After (F1T41.22), the fixed-weight space of
\(D_\tau\) is \(E_\tau^*\otimes\mathcal H_0\).
Its large stabilizer is irreducible: an invariant
orthogonal projection commuting with every matrix
unit on \(\mathcal H_0\) is \(T\otimes I\),
and commuting with the right-regular constituent
forces \(T\) scalar. Thus \(D_\tau\) has exactly
one top \(H\)-constituent \(\eta_\tau\), once;
every other constituent is in its lower part
\(B_\tau\). Evaluation at a nonzero vector of
\(E_\tau\) injects its actual multiplicity space
into \(\mathcal R_k\), so \(\dim B_\tau\le E\).
All selected original minimal-level constituents
are among these \(\eta_\tau\).

The averaged projector on the actual smaller
action, divided by \(b\), is exactly
\[
 \operatorname{Tr}_{D_\tau}(x)
       ={1\over s}\sum_{t\in S}\overline{\tau(t)}
                               \operatorname{Tr}\mathcal R_k(x,t).
 \tag{F1T41.23}
\]
The anisotropic coefficient factor has been
retained in both the irreducibility proof and
this normalization.

\paragraph{7. The literal commutator bound and the original source contradiction.}
The translation/multiplication basis of
\(\operatorname{End}\mathbf C[A^k]\) is labelled
by \(A\otimes U\), using the nondegenerate
polar form \(B_A\). The actual \(H\times S\)
operators conjugate its labels by \(x\otimes t\),
with their unit phases retained. Directly on
the three smaller generators this follows from
the column, quadratic-phase and Fourier
formulas in paragraph 3. Thus
\[
       |\operatorname{Tr}\mathcal R_k(x,t)|^2
                   \le q^{\dim\ker(x\otimes t-I)}.
 \tag{F1T41.24}
\]
For \(t\ne I\), the partial Fourier model gives
\[
 |\operatorname{Tr}\mathcal R_k(I,t)|
   =q^{r\dim\ker(t-I)}|\operatorname{Tr}\rho_0(t)|
                             \le Dq^{-r}.
 \tag{F1T41.25}
\]
For the same literal \(h\) in (F1T41.6),
\[
 |\operatorname{Tr}\mathcal R_k(h,I)|\le Dq^{-2k},
 \qquad
 |\operatorname{Tr}\mathcal R_k(h,t)|
             \le Dq^{-(N-4)/2}\le Dq^{-r+2}\quad(t\ne I).
 \tag{F1T41.26}
\]
For the cross term, restrict \(h\otimes t-I\)
to \(\ker(h-I)\otimes U\); its rank is at least
\((N-4)\operatorname{rank}(t-I)\).

Set \(B=bD/s\). Equations (F1T41.17) and
(F1T41.23)--(F1T41.26) give
\[
\begin{split}
 \dim\eta_\tau&>
 B\bigl(1-q^{4k^2-r}-q^{4k^2+2k-N}\bigr),\\
 |\operatorname{Tr}\eta_\tau(h)|&<
 B\bigl(q^{-2k}+q^{4k^2-r+2}+q^{4k^2+2k-N}\bigr).
\end{split}
 \tag{F1T41.27}
\]
Since \(N\ge2r\) and \(r\ge4k^2+9\), the
same exact endpoint as in T40 gives
\[
          { |\operatorname{Tr}\eta_\tau(h)|\over\dim\eta_\tau}
                          <{67585\over261119}<{7\over27}.
 \tag{F1T41.28}
\]
Summing over the actual original restriction
and every multiplicity gives
\( |\operatorname{Tr}_V\beta(h)|<7m/27\).
The original finite-group Hermitian form gives
\(\operatorname{rank}(\beta(h)-I)>10m/27\), and
the literal commutator identity bounds this by
\(2\operatorname{rank}(g-I)\).
Thus each nonscalar original inertia has moving
rank greater than \(5m/27\); a nonidentity scalar
has moving rank \(m\). The original budget
\(\sum_{j=1}^b\operatorname{rank}(g_j-I)=2m+6\)
with \(b\ge12\) forces \(m<27\), contradicting
the retained source gate. This proves (F1T41.1).

\paragraph{8. Exact framework and source scope.}
The displayed form-isometry conjugations transport
the natural cover map by its explicit inverse.
All other comparisons are the actual Hom
evaluation, Fourier change of basis, genuine
small-group restriction/inclusion/evaluation,
large unipotent action, code embedding, full-weight
unitary basis change and averaged projector.
They operate on constituents of the original
coefficient representation. The original curve,
parabolic cohomology and primitive Jacobian corner
are unchanged. Their cup, restricted polarization,
Poincar\'e biextension, period and Grothendieck--Adams
maps therefore remain the identity, with pairing
multiplier one. The bound adds no detector family.

The named primary inputs are the actual T35--T36
cover/level statements, GLT Lemma 3.3 with its
recorded parity dependencies, Steinberg Theorem 30's
finite-field proof and Theorem 36's supplied sketch,
the actual odd Weil linearization proved in T37,
and the genuine unitary-Heisenberg extension with its
published proof scopes audited in T39. All new anisotropic code and stabilizer
claims are proved by the displayed operators above.
The accompanying source audit records precise
fresh and retained scopes, including the recursive
limits of the named classification inputs.
